\section{Introduction} \label{sec:intro}

Modern systems engineering is increasingly shaped by the mandate to adopt a Modular Open Systems Approach (MOSA) in the acquisition, architecture, and sustainment of complex cyber-physical systems~\cite{esper2019modular, zimmerman2019considerations}. Long-range digital-engineering visions similarly emphasize composable, model-based practice across the system lifecycle~\cite{INCOSE2035}. Reference architectures such as the Future Airborne Capability Environment (FACE)\textsuperscript{\textregistered} Technical Standard~\cite{OpenGroup2023FACE}, the Sensor Open Systems Architecture (SOSA)\textsuperscript{\textregistered} Reference Architecture~\cite{SOSAConsortium2023}, and Open Mission Systems / Universal Command and Control Interface (OMS/UCI)~\cite{esper2019modular, tokar2017osa} seek to replace monolithic, vendor-locked baselines with modular ecosystems in which hardware and software capabilities can be independently developed, inserted, and upgraded. Yet operational experience across aerospace, defense, and autonomous systems demonstrates that standardized transport buses and open interface boundaries do not, by themselves, deliver rapid or automated systems integration~\cite{madni2014system}. Even when two subsystems conform to the same reference standard, integrating an upstream \emph{Producer} capability with a downstream \emph{Consumer} capability routinely demands substantial manual engineering effort to reconcile discrepancies in domain meaning, coordinate reference frames, units of measure, wire-level type encodings, and asynchronous interaction timing.

At its core, this integration bottleneck reflects a structural division in how current model-based systems engineering (MBSE) methodologies represent interfaces. Existing formalisms largely separate \emph{static data semantics} from \emph{dynamic system behavior} and \emph{architectural trade-offs}. On the semantic side, standards such as the Open Universal Domain Description Language (UDDL)\textsuperscript{\textregistered}~\cite{OpenGroup2023UDDL} separate domain meaning from measurement convention and machine representation across conceptual, logical, and platform layers, enabling ontology alignment and query-translation algorithms to identify static correspondences between heterogeneous data models~\cite{mckibben2026alignment, mckibben2026nearring}. Identifying that two interface elements refer to related domain concepts, however, does not synthesize the stateful mediation architecture needed to couple them dynamically—such as buffering a high-rate pushed sensor stream for an on-demand guidance consumer, correlating asynchronous attitude and position reports to rotate reference frames, or bounding numerical precision loss and staleness under real-time execution constraints. Conversely, formal systems-theoretic frameworks—most notably A.~Wayne Wymore's Tricotyledon Theory of System Design (T3SD)~\cite{Wymore1993}—provide a rigorous algebra for coupling discrete component systems from a technology catalog to satisfy trajectory-level input/output requirements, but treat component ports as unstructured set-theoretic alphabets devoid of domain semantics.

Viewed through a systems-theoretic lens, a MOSA component ecosystem naturally mirrors Wymore's buildability formulation: the set of available Producer capabilities, Consumer capabilities, and reusable Mediation Operators constitutes a \emph{Technology Catalog}, and a concrete integration architecture is a \emph{System Coupling Recipe} whose synchronous composition yields a closed-loop \emph{Coupled Resultant}. In classical T3SD, certifying that a coupled assembly satisfies a behavioral specification requires manually designing the coupling recipe and constructing a surjective system homomorphism onto an abstract functional model—a manual effort that does not scale to combinatorial integration spaces or evolving component catalogs. If the port compatibility constraints of the Technology Catalog are enriched with multi-tier UDDL Semantic Colors and Flow Kinematics, however, the synthesis of a valid mediation architecture—and the online regulation of its runtime traffic—can be cast as a well-defined sequential decision problem over a formal design space.

This paper investigates whether framing semantically enabled MOSA integration as a formal game over a semantics-refined Wymorian buildability space provides a rigorous foundation for automated verification and learning-guided synthesis. Specifically, we organize our inquiry around three research questions:
\begin{enumerate}[label=\textbf{RQ\arabic*:}, leftmargin=2.8em, topsep=3pt, itemsep=2pt]
    \item \emph{Formalization of Semantic-Behavioral Integration:} How can system integration under a Modular Open Systems Approach be formalized when Producer Offers and Consumer Needs diverge simultaneously across conceptual entities, logical measurement frames, platform encodings, and interaction flow kinematics?
    \item \emph{Trajectory-Level Behavioral Satisfiability:} What formal conditions ensure that a synthesized mediation architecture not only aligns static interface types and flow mechanics, but also satisfies dynamic, trajectory-level behavioral contracts without overfitting finite test scenarios?
    \item \emph{Learning-Guided Architecture Synthesis and Control:} How can reinforcement learning leverage multi-tier semantic structure, axiomatic cost trade-offs, and assertional behavioral monitors to navigate the combinatorial space of coupling recipes offline and regulate fault-tolerant mediation policies online?
\end{enumerate}

To address these questions, we introduce \emph{The Integration Game}, a systems-theoretic methodology for dynamic interface mediation. Rather than treating interface compatibility as a static documentation checklist, we synthesize UDDL's three-tier data semantics, Wymore's T3SD coupling algebra, and Universal Live Sequence Charts (uLSCs)~\cite{DammHarel2001} into an executable design framework. Each atomic \emph{Field} exposed at a component boundary is typed by a three-axis \emph{Semantic Color} $c = (C, L, P)$ together with its \emph{Flow Kinematics} $\mathbf{k} = (x, v, a)$ and payload metadata. Legal moves on the \emph{Mediation Board} correspond to primitive \emph{Mediation Operators} that recolor Fields, aggregate or decompose Message structures, derive fused observables, buffer state, and bridge asynchronous interaction modes. To evaluate behavioral correctness without manual homomorphism proofs, we compile Wymorian Input/Output Requirements ($IOR$s) into assertional uLSCs and prove their behavioral equivalence. We then formulate a generalized \emph{Integration Cost Functional} governed by seven technology-independent axioms, and show how compiling the uLSC into a finite-state Reward Machine~\cite{ToroIcarte2022} paired with Counterexample-Guided Inductive Synthesis (CEGIS) enables Reinforcement Learning (RL) agents to synthesize verified mediation architectures offline and adapt routing policies online.

The remainder of the paper is organized as follows. Section~\ref{sec:background} establishes the theoretical foundations of UDDL data semantics, Wymore's T3SD cotyledons, and Universal Live Sequence Charts, and situates our contribution relative to prior work. Section~\ref{sec:methodology} develops the formal model of \emph{The Integration Game}, including typed Wymorian coupling, the taxonomy of Mediation Operators, the Three-Tier Winning Condition, the axiomatic Integration Cost Functional, and the Reward Machine / CEGIS learning formulation. Section~\ref{sec:results} evaluates the methodology across a fifteen-scenario progressive benchmark suite, reporting cost-sensitivity audits, symbolic fragment verification, offline synthesis generalization, and online fault-recovery performance. Section~\ref{sec:discussion} examines the systems-engineering implications, decidability boundaries, and methodological scope of the approach, and Section~\ref{sec:conclusion} concludes.
